\chapter{Flight Software}

\section{Mission Phases}

The mission is divided into six flight phases, each representing a physically distinct regime of the stratospheric balloon flight.

\begin{table}[H]
\centering
\begin{tabular}{lp{10cm}}
\toprule
\textbf{Phase} & \textbf{Physical description} \\
\midrule
Standby & System powered on, sitting on ground. All sensors initialising. Ground pressure baseline (P\_ground) recorded here as reference for all altitude calculations. \\
Launch & Balloon released. Sharp upward jerk, brief chaotic motion lasting seconds. Transition from ground-constrained to free flight. \\
Ascent & Balloon stabilised, climbing at 3-6 m/s steadily. Longest active phase, typically 1.5-2 hours. Temperature drops significantly with altitude. \\
Cruise & Balloon fully expanded at float altitude (~25-35 km). Vertical velocity near zero, altitude oscillating within a bounded range. Duration 30 min to several hours. \\
Descent & Balloon envelope bursts. Rapid fall begins, IMU sees sudden shift from near 0g float to ~1g. Parachute deploys, stabilising descent rate. \\
Landing & System touches down. All motion stops. Terminal state, no further FSM transitions. Logging and telemetry continue for recovery; shutdown timer starts. \\
\bottomrule
\end{tabular}
\caption{Mission phases and physical description}
\end{table}

\section{Complete Finite State Machine}
\label{sec:fsw-fsm}

The FSM is a bare-metal Moore machine, implemented as a single control
loop running at 1~Hz on the STM32L476RG. It consists of six states and
five transition paths, including one direct path from Ascent to Descent
used if the balloon bursts before float is reached.

\subsection{State Transition Table}

\begin{table}[H]
\centering
\begin{tabular}{p{2.1cm}p{2.0cm}p{9.6cm}}
\toprule
\textbf{Current state} & \textbf{Next state} & \textbf{Transition condition} \\
\midrule
Standby & Launch & IMU total acceleration magnitude $>$1.5g (median
filtered, 2 consecutive seconds) AND barometric altitude increase
$>$10~m above $P_\text{ground}$ baseline over the same window. \\
Launch & Ascent & Barometric altitude $>$100~m above $P_\text{ground}$
baseline AND GPS vertical velocity $>$2~m/s upward (fallback:
barometric-derived vertical velocity if GPS fix invalid), sustained
3~s. \\
Ascent & Cruise & GPS vertical velocity $<$0.5~m/s (fallback:
barometric-derived) AND barometric altitude bounded within
$\pm$50~m over window, sustained 30~s. \\
Ascent & Descent & Direct path if Cruise never detected (balloon burst
before float). Same condition as Cruise to Descent below, applied
directly from Ascent. \\
Cruise & Descent & IMU total acceleration deviates $>$0.3g from running
float median AND GPS vertical velocity $<-5$~m/s (fallback:
barometric-derived), sustained 5~s. \\
Descent & Landing & GPS total 3D speed $<$2~m/s (fallback: GPS vertical
velocity, then barometric-derived) AND barometric altitude $<$50~m
above $P_\text{ground}$ baseline, sustained 5~s. \\
Landing & Terminal & No further transitions. Shutdown timer starts on
entry; non-essential systems powered down after a defined period. \\
\bottomrule
\end{tabular}
\caption{FSW finite state machine: states, transitions, and conditions}
\label{tab:fsw-transitions}
\end{table}

\section{Sensor Parameter Reference}
\label{sec:fsw-sensors}

Three sensors are used for phase detection: MPU-6050 (IMU), u-blox
MAX-M10S (GPS), MS5607 (barometer). Each provides direct measurements
and, where relevant, derived parameters obtained through
differentiation.

\subsection{IMU (MPU-6050)}

\begin{table}[H]
\centering
\begin{tabular}{p{3.2cm}p{2.3cm}p{2.1cm}p{6.4cm}}
\toprule
\textbf{Parameter} & \textbf{Type} & \textbf{Reliability} & \textbf{Notes} \\
\midrule
Acceleration X/Y/Z & Direct & High & Reads $\sim$1g at rest on the
Z axis (normal force), not 0. Cannot distinguish steady ascent from
stationary. \\
Total acceleration magnitude & Derived (1 step) & High &
Orientation-independent event detection. Used for Standby to Launch and
Cruise to Descent burst. \\
Angular rate X/Y/Z & Direct & High & Near zero at rest and during smooth
flight. Not useful for altitude-based phases. \\
Velocity (integrated) & Derived (1 step) & Low & Drift accumulates
rapidly; only reliable over a few seconds. Not used for phase
detection. \\
Position (double integrated) & Derived (2 steps) & Very low & Not used;
drift makes this unreliable beyond seconds. \\
\bottomrule
\end{tabular}
\caption{IMU parameters}
\label{tab:fsw-imu}
\end{table}

\subsection{GPS (u-blox MAX-M10S)}

\begin{table}[H]
\centering
\begin{tabular}{p{3.2cm}p{2.3cm}p{2.1cm}p{6.4cm}}
\toprule
\textbf{Parameter} & \textbf{Type} & \textbf{Reliability} & \textbf{Notes} \\
\midrule
Altitude (WGS84) & Direct & Medium & Vertical accuracy $\sim$5--15~m
error, worse than horizontal. Unreliable if no fix. \\
Vertical velocity & Direct & High & Doppler-based direct measurement.
Requires valid fix; cold-start delay at power-on; COCOM limit above
18~km if not in Airborne mode. \\
Total 3D speed & Derived (1 step) & High & Primary condition for
Descent to Landing; catches horizontal drift as well as vertical. \\
Fix status & Direct & High & Gate condition: all other GPS parameters
ignored if fix invalid, fallback to barometer applies. \\
GPS-derived acceleration & Derived (1 step) & Medium & Cross-check
against IMU acceleration for fault detection. \\
\bottomrule
\end{tabular}
\caption{GPS parameters}
\label{tab:fsw-gps}
\end{table}

\subsection{Barometer (MS5607)}

\begin{table}[H]
\centering
\begin{tabular}{p{3.2cm}p{2.3cm}p{2.1cm}p{6.4cm}}
\toprule
\textbf{Parameter} & \textbf{Type} & \textbf{Reliability} & \textbf{Notes} \\
\midrule
Pressure & Direct & High & Raw input, not used directly in phase
logic. Requires temperature compensation using factory calibration
coefficients. \\
Altitude (relative to $P_\text{ground}$) & Derived (1 step) & High &
Primary altitude source for most transitions. Valid across the full
0--30$+$~km range. \\
Vertical velocity (derived) & Derived (2 steps) & Medium & Noisier than
GPS direct; used as a fallback confirming condition when GPS fix is
invalid. \\
Vertical acceleration (derived) & Derived (2 steps) & Low & Not used as
primary; too noisy for sharp event detection. \\
\bottomrule
\end{tabular}
\caption{Barometer parameters}
\label{tab:fsw-baro}
\end{table}

\section{Threshold Values and Justification}
\label{sec:fsw-thresholds}

All threshold values are set conservatively relative to observed ranges
from real stratospheric balloon flight telemetry where available, or
from engineering margin based on sensor characteristics where a specific
citable figure was not found.

\begin{table}[H]
\centering
\small
\begin{tabular}{p{2.1cm}p{2.7cm}p{1.9cm}p{7.0cm}}
\toprule
\textbf{Transition} & \textbf{Condition} & \textbf{Threshold} & \textbf{Justification} \\
\midrule
Standby to Launch & IMU acceleration magnitude (primary) & $>$1.5g, 2~s
& Resting value $\sim$1.026g \cite{safonova2016}. Threshold set as
engineering margin above resting noise to reliably capture the release
jerk. \\
Standby to Launch & Barometric altitude change (confirm) & $>$10~m &
MS5607 altitude resolution 20~cm per datasheet \cite{ms5607}. Ground
level noise typically 1--2~m; 10~m provides clear margin above this. \\
Launch to Ascent & Barometric altitude (primary) & $>$100~m & Real
ascent rates of 4.6--5.9~m/s \cite{safonova2016, winterballoon2018}
reach 100~m within 20--25~s of release. \\
Launch to Ascent & GPS vertical velocity (confirm) & $>$2~m/s, 3~s & Set
conservatively below the observed 4.6--5.9~m/s ascent rate range
\cite{safonova2016, winterballoon2018}, ensuring the transition fires
even for a slower-filling balloon. \\
Ascent to Cruise & GPS vertical velocity (primary) & $<$0.5~m/s, 30~s &
A sharp stop at float altitude of 33{,}200~m with no overshoot was
recorded, confirming the ascent-to-float transition is abrupt rather
than gradual \cite{winterballoon2018}. 0.5~m/s set as engineering margin
below this. \\
Ascent to Cruise & Barometric altitude bounded (confirm) & $\pm$50~m &
Engineering margin; no specific citable oscillation figure identified
for this exact bound. \\
Cruise to Descent & IMU deviation from float baseline (primary) &
$>$0.3g, 5~s & Float acceleration is stable and low-noise, average
1.026g with low RMS motion \cite{safonova2016}. Threshold set as
engineering margin above this stable baseline to detect burst onset. \\
Cruise to Descent & GPS vertical velocity (confirm) & $<-5$~m/s & Real
burst events produce descent rates up to $-52$~m/s immediately after
burst, stabilising around $-10$~m/s under parachute \cite{oh3vhh2021}.
$-5$~m/s set conservatively above the stable descent rate. \\
Descent to Landing & GPS total 3D speed (primary) & $<$2~m/s, 5~s &
Terminal velocity under parachute of 5.5~m/s confirmed \cite{cusfnova1}.
2~m/s set as engineering margin well below this. \\
Descent to Landing & Barometric altitude (confirm) & $<$50~m & MS5607
resolution of 20~cm per datasheet \cite{ms5607} makes 50~m an
unambiguous near-ground indicator. \\
\bottomrule
\end{tabular}
\caption{Threshold values and justification}
\label{tab:fsw-thresholds}
\end{table}

\section{FSM Classification: Moore Machine}
\label{sec:fsw-moore}

The FSW FSM is classified as a Moore machine.

FSW's outputs and behaviours -- baseline recording, watchdog servicing,
GPS Airborne-mode monitoring, IMU float-baseline tracking, shutdown
timer -- are determined solely by the current state, not by the specific
instantaneous sensor input value. Sensor readings are used only as
inputs to evaluate transition conditions; they do not themselves alter
the output behaviour while the system remains in a given state. For
example, whether the IMU reads 1.02g or 1.04g while in Ascent, the
system's actions (poll, filter, check watchdog) remain identical; only
whether a transition condition is satisfied changes. This differs from a
Mealy machine, where outputs would depend on the specific input value in
addition to the state. Logging rate and telemetry rate, owned by CDH and
TT\&C respectively, are keyed to current phase rather than instantaneous
sensor values, consistent with this classification.
